\section{Conclusions}
\label{sec:conclusion}

We have developed a physics-based model for the acoustic `thump test''
used to assess watermelon ripeness.  The principal findings are:

\begin{enumerate}
  \item The watermelon is well described as a fluid-filled viscoelastic
    oblate spheroidal shell.  The dominant tap-tone corresponds to the
    $n = 2$ flexural mode, with frequency~$f_2$ given by
    Eq.~\eqref{eq:fn_general}.

  \item The squared frequency $f_2^2$ is exactly linear in the rind
    elastic modulus $E_{\mathrm{rind}}$ when turgor pressure is
    negligible.  This linearity enables a closed-form inversion
    (Eq.~\ref{eq:inversion}) from a single tap-tone measurement to the
    rind modulus.

  \item Sobol sensitivity analysis confirms that
    $E_{\mathrm{rind}}$ dominates the frequency response
    ($S_T > 0.5$), with fruit size as the second most influential
    parameter.

  \item The dimensionless ripeness parameter $\Pi_{\mathrm{ripe}}$
    (Eq.~\ref{eq:Pi_ripe}) collapses predictions across cultivars onto
    a universal master curve, providing a cultivar-independent ripeness
    metric.

  \item Model predictions agree with published resonance measurements
    on watermelons \cite{Yamamoto1980} and other fruits
    \cite{DeBelieetal2000, ChenDeBaerdemaeker1993}.
\end{enumerate}

The framework generalises directly from prior work on biological
cavities, confirming that classical shell theory provides a unified
language for fluid-filled biological structures---whether the application
is clinical diagnostics or selecting the perfect watermelon.

Future work will pursue experimental validation on instrumented
watermelons, extension to other fruits (cantaloupe, coconut, pumpkin),
and development of a smartphone-based measurement tool.
